\chapter*{Introduction}
\addcontentsline{toc}{chapter}{Introduction}

Le bloc opératoire concentre des ressources coûteuses, rares et fortement interdépendantes. La programmation d'une intervention mobilise non seulement une salle et une vacation opératoire, mais aussi un chirurgien, une équipe d'anesthésie, du personnel soignant et, selon le parcours du patient, une place ambulatoire ou un lit d'hospitalisation. Une décision prise sur une seule de ces ressources peut donc déplacer la contrainte vers une autre partie de l'établissement : maximiser l'occupation des salles peut saturer les lits, tandis qu'un pilotage centré uniquement sur l'hébergement peut dégrader l'utilisation du bloc. Les revues récentes recommandent ainsi une vision systémique et une planification intégrée des ressources plutôt qu'une optimisation locale du seul bloc opératoire \parencite{rachuba2024,schouten2023}.

Dans une organisation en \emph{flux poussé}, la date est d'abord choisie à partir du planning du chirurgien, puis les autres capacités doivent absorber cette décision. Cette logique contribue à des alternances de pics et de creux d'activité, particulièrement difficiles à anticiper autour des week-ends et des périodes de vacances. Le projet étudie une logique de \emph{flux tiré} : à partir des caractéristiques du patient et des capacités prévisionnelles, un outil d'aide à la décision propose une ou plusieurs dates compatibles. Le professionnel conserve la décision finale.

La problématique étudiée est ainsi la suivante : \emph{comment exploiter les données historiques de séjour et d'intervention pour prévoir les besoins d'un patient, puis construire un ordonnancement qui utilise efficacement le bloc tout en lissant l'occupation des lits et en restant robuste aux aléas ?}

Pour répondre à cette question, le travail poursuit quatre objectifs. Le premier est de comprendre la base fournie, d'en mesurer la qualité et de construire des variables opérationnelles fiables. Le deuxième est de prédire des quantités utiles avant l'admission, notamment la durée de séjour et le temps d'occupation d'une salle. Le troisième est de formaliser les tâches, les ressources, les contraintes et les critères qui définissent un planning réalisable. Le quatrième est d'implémenter puis de comparer plusieurs métaheuristiques, en particulier la recherche taboue, le recuit simulé et les algorithmes génétiques.

Le rapport présente d'abord le problème d'ordonnancement des blocs et ses principales difficultés. Il analyse ensuite la base de données et prépare la démarche prédictive. Les chapitres suivants détaillent l'adaptation et l'évaluation des méthodes d'optimisation sur des instances de tailles différentes, avant une comparaison globale et une discussion des perspectives. Les recommandations produites devront rester explicables, être validées dans le contexte réel de l'établissement et compléter, sans la remplacer, l'expertise des professionnels de santé.
